\chapter{Applications Beyond Autonomous Vehicles}
\label{sec:applications}
% ============================================

% Springer Nature: each chapter should have a 10--15 line online abstract.
% This chapter abstract has not yet been written in the uploaded manuscript.
\abstract*{}
\suppressfloats[t]

%\deleted{[This paragraph doesn't make much sense and anyway not here. Seems a leftover]} % Resolved: the leftover was replaced by the bridging paragraph below (blue).
\added{The two case studies in this section are deliberately chosen to show how \emph{all} the classical ingredients of Chapter~\ref{sec:classicalSF} --- Kalman filtering and innovation monitoring, evidence-theoretic conflict handling, and MCDM-based prioritisation --- combine into complete fusion systems in a domain far from autonomous driving. They should be read as integrative worked examples of the survey's toolbox rather than as an exhaustive treatment of maritime sensing.}

\section{Maritime Anomaly Detection and Infrastructure Protection}

Maritime anomaly detection is a representative application in which sensor fusion is needed because no single information source provides a complete and trustworthy maritime picture. The Automatic Identification System (AIS) provides vessel identity, position, speed, course, and other navigation information, but AIS messages can be missing, delayed, spoofed, or intentionally misleading. Coastal radar, electro-optical and infrared cameras, satellite imagery, synthetic-aperture radar, hydrophones, and distributed acoustic sensing can provide independent observations, but each modality has different coverage, resolution, latency, and false-alarm characteristics. Recent surveys of AIS anomaly detection highlight the diversity of abnormal behaviours that must be considered, including route deviation, loitering, unexpected stops, speed anomalies, encounter anomalies, and suspicious gaps in transmission \cite{wolsing_anomaly_2022}. This problem is becoming more relevant as maritime navigation systems become increasingly digital and interconnected, exposing the fused navigation pipeline itself to faults, spoofing, and cyber--physical disturbances \cite{oliva_innovation_2026}.

The practical difficulty is that suspicious behaviour is not always dramatic. A vessel may keep transmitting AIS in order to avoid immediate suspicion while still performing small route deviations, abnormal speed profiles, unusual loitering, or manoeuvres close to protected infrastructure. This is especially relevant in congested and geopolitically sensitive regions such as the Baltic Sea, where normal maritime traffic coexists with dense networks of submarine power and communication cables. A maritime fusion architecture should therefore combine trajectory analysis with contextual information: proximity to coastlines, shipping lanes, restricted areas, cable corridors, anchorages, recent incidents, weather, and vessel class.

At the trajectory level, fusion typically starts by constructing vessel tracks from AIS and radar measurements, associating visual or satellite detections with those tracks, and extracting behavioural features such as speed over ground, course change, distance from shipping lanes, proximity to restricted areas, encounter geometry, and consistency between reported identity and observed motion. Remote-sensing imagery and AIS can also be fused to identify small vessels or vessels whose reported AIS state is inconsistent with the observed scene \cite{chen_sviadf_2025}. Graph-based route models provide an additional layer of context: a recent Baltic Sea framework constructs a Probabilistic Roadmap (PRM) of feasible maritime corridors, projects AIS trajectories onto this graph, and derives interpretable features such as route inefficiency, deviation from learned corridors, and speed or dwell time near submarine cables \cite{oliva_tomei_prm_access_2026}. This illustrates a broader design principle: raw trajectories become more informative when fused with geospatial constraints and infrastructure maps.

Maritime anomaly detection should also include the estimation layer. In navigation and monitoring systems, GNSS, inertial sensors, radar, and AIS may be fused through Kalman-filter-based estimators. Some cyber--physical anomalies or stealthy disturbances may first appear as abnormal innovations before they become visible as large trajectory deviations. Innovation-based monitoring can therefore test whether the residual sequence remains zero-mean, covariance-consistent, Gaussian, and temporally white; such diagnostics can expose coloured spoofing or disturbance patterns that may bypass a simple \(\chi^2\) covariance check \cite{oliva_innovation_2026}. This connects ship anomaly detection directly to the sensor-fusion pipeline rather than treating it only as offline post-processing of AIS tracks.

Evidence theory is well suited to this setting because maritime anomalies are often ambiguous. A vessel moving slowly near a cable route may be fishing, drifting because of a mechanical problem, performing maintenance, or preparing an illegal anchoring activity. A frame of discernment can include hypotheses such as \emph{normal transit}, \emph{loitering}, \emph{route deviation}, \emph{AIS inconsistency}, \emph{collision risk}, \emph{infrastructure threat}, and \emph{unknown}. AIS-derived behaviour, radar tracks, camera detections, satellite observations, innovation-test alarms, and acoustic or cable-sensing cues can each provide BBAs over this frame. If several sources support the same abnormal hypothesis, the fused belief increases. If AIS reports normal transit while radar, imagery, or cable sensing indicates inconsistent behaviour, the conflict mass becomes an operator-relevant warning rather than being hidden inside a single confidence score.

MCDM can complement this process by ranking alerts and response actions. Criteria may include distance to protected infrastructure, confidence of track association, deviation from learned maritime corridors, innovation-test alarms, severity of the suspected anomaly, persistence over time, weather and visibility, vessel type, traffic density, and expected operator workload. AHP-style pairwise comparisons can encode expert preferences from coast guards, port authorities, or infrastructure operators, while real-time sensor indicators provide the criterion scores. The resulting MCDM score can prioritise which anomalies require immediate attention, which sensors should be tasked for confirmation, or which vessels should be monitored more closely.

The main challenges are data quality, class imbalance, interpretability, and validation. True anomalous events are rare and often poorly labelled, while normal maritime traffic varies strongly with geography, season, weather, and port operations. Synthetic anomalies, incident-based case studies, and infrastructure-aware features can help evaluate candidate methods, but broader external validation remains necessary \cite{oliva_tomei_prm_access_2026}. For operational deployment, a transparent combination of track-level fusion, innovation monitoring, evidence-theoretic uncertainty handling, and MCDM-based prioritisation can be more useful than a black-box anomaly score alone.

\section{Sensor Fusion for Submarine Cable Monitoring}

Submarine power and telecommunication cables are critical infrastructures for digital connectivity, energy interconnection, offshore wind farms, and security-sensitive maritime operations \cite{carter_submarine_2009,kavanagh_achieving_2025}. They are exposed to accidental and intentional hazards such as anchor dragging, bottom trawling, seabed instability, fatigue, construction activity, and deliberate tampering. Anchor-related events are particularly important, with the risk depending on anchor geometry, burial depth, seabed properties, and vessel drift \cite{yoon_safety_2013,osthoff_anchor_soil_2017,li_anchor_risk_2025}. Monitoring is difficult because the relevant events are spatially distributed, the assets extend over tens or hundreds of kilometres, and damaging interactions may be masked by waves, currents, seabed coupling, and ordinary vessel activity.

Distributed acoustic sensing (DAS) turns an optical fibre into a dense array of strain or strain-rate measurements by interrogating backscattered light along the cable \cite{hartog_intro_2017,lu_distributed_2019,lindsey_fiberoptic_2021}. In marine environments, this is attractive because seafloor instrumentation is difficult and expensive to deploy, while existing or dedicated fibre-optic cables may cover long distances with metre-scale spatial sampling. Long-term subsea deployments show that DAS can support real-time detection of bottom-trawl fishing, ship anchoring, and seismic events over cable spans of operational length \cite{waagaard_longterm_2022}, and related demonstrations cover seismic, hydroacoustic, and power-cable strain sensing \cite{williams_distributed_2019,matsumoto_hydroacoustic_2021,masoudi_dynamic_2018}.

For underwater-cable monitoring, DAS can support several fusion tasks. First, it can provide spatially distributed acoustic or deformation evidence of nearby vessels, anchors, trawling activity, sediment motion, or cable vibration. Second, it can be fused with AIS, radar, bathymetry, weather, ocean-current estimates, and maintenance records to interpret the cause of the observed signal. Third, it can support infrastructure protection by detecting events close to telecom cables, power cables, pipelines, offshore wind farms, or port approaches. Experimental work using an optical telecom cable offshore Toulon showed that ship-radiated noise could be detected and, in favourable conditions, used for trajectory estimation \cite{rivet_ship_2021}, while more recent Applied Ocean Research work studies DAS-based vessel monitoring using submarine fibre-optic cables in the presence of growing offshore infrastructure and vessel traffic \cite{paap_vessel_2025}. For power cables, \(\phi\)-OTDR experiments have also been used to distinguish calm sea, anchor drop, and anchor drag signatures \cite{lv_otdr_2018}.

The fusion problem is challenging because DAS measurements are not equivalent to conventional hydrophone measurements. The response depends on cable coupling to the seabed, cable orientation relative to the incoming wavefield, gauge length, interrogation settings, local bathymetry, and the propagation path. Some cable segments may be well coupled to the seafloor, while suspended or poorly coupled sections may respond differently. A robust fusion system should therefore include a cable-state model and should not treat all DAS channels as equally reliable.

Model-based monitoring can help bridge the gap between raw distributed measurements and actionable alarms. Cable dynamics can be approximated through physically grounded tensioned-structure models, with finite-dimensional state-space approximations providing a natural interface to filtering and diagnosis \cite{irvine_linear_1974,baudouin_cable_pde_2018}. 

In a recent paper~\cite{fioravanti_cable_med_2026} this direction is followed by modelling a submarine cable as a damped tensioned structure, obtaining a finite-dimensional state-space approximation, and applying Kalman filtering to distributed deformation measurements. Under nominal conditions, the filter innovations should remain statistically consistent with the model and the assumed environmental excitation. Localised mechanical interactions, such as an anchor strike or tampering event, appear as systematic deviations in the innovation sequence and can be detected using sliding-window statistical tests. This approach complements purely data-driven DAS processing because it introduces physical interpretability and separates stochastic environmental forcing from persistent structural perturbations.

Evidence theory can be used to represent uncertainty in event interpretation. For a cable segment and time window, the hypotheses may include \emph{passing vessel}, \emph{anchoring or trawling risk}, \emph{seismic event}, \emph{oceanographic background}, \emph{cable-coupling artefact}, \emph{mechanical tampering}, and \emph{unknown}. DAS spectral features, distributed deformation residuals, AIS proximity, radar tracks, hydrophone data, bathymetry, and environmental models can each assign belief masses. High conflict, such as a strong cable disturbance without a nearby AIS vessel, may indicate an unreported vessel, missing data, a non-vessel geophysical source, or a localised mechanical interaction requiring inspection.

MCDM is useful for prioritising alerts and sensor resources. Criteria may include signal-to-noise ratio, distance to the cable, estimated vessel size, consistency with AIS, innovation-test significance, persistence of the event, cable criticality, environmental conditions, and cost of dispatching inspection assets. The same framework can rank candidate explanations or decide whether to request additional confirmation from patrol vessels, satellite imagery, or nearby fixed sensors. In this way, DAS-based cable monitoring becomes a full sensor-fusion problem: dense acoustic or deformation sensing provides the raw evidence, physical models and filters structure the residuals, and contextual maritime data and decision logic turn that evidence into actionable situation awareness.
